\section{并行工作与多 Agent 协作}
%% ============================================================

Claude Code 支持多种并行工作模式，从单机 worktree 到多 Agent 团队协作。

\subsection{用 --worktree 做并行隔离开发（Tip \#15）}

\texttt{claude --worktree feature-auth} 创建一个带新分支的隔离工作副本。Claude Code 团队称这是最大的生产力解锁之一。

\begin{importantbox}{Worktree 最佳实践}
通常同时运行 2--3 个 worktree，每个有独立会话、独立分支、独立文件系统状态。本地 worktree 的上限是你的机器资源（CPU、内存）。如需更多并行度，可以考虑云容器方案（如 Builder.io）。
\end{importantbox}

\subsection{用 Subagent 保持主上下文干净（Tip \#19）}

"Use subagents to figure out how the payment flow handles failed transactions." 这会生成一个独立的 Claude 实例，有自己的上下文窗口。它读取文件、推理后返回简洁摘要。

内置类型：
\begin{itemize}
  \item \textbf{Explore}（Haiku）：快速文件搜索
  \item \textbf{Plan}（只读分析）：代码库推理
\end{itemize}

深度调查可能消耗一半上下文窗口，Subagent 将这些成本隔离在主会话之外。

\subsection{自定义 Subagent（Tip \#35）}

与临时使用不同，自定义 Subagent 保存在 \texttt{.claude/agents/}，预配置好模型和工具权限。例如：用 Opus 和只读工具的 security-reviewer agent，或用 Haiku 追求速度的 quick-search agent。

用 \texttt{/agents} 浏览和创建。可以设置 \texttt{isolation: worktree} 让 agent 拥有独立文件系统。

\subsection{Agent Teams：多会话协调（Tip \#20）}

\begin{knowledgebox}{实验性功能}
需要先启用 \texttt{CLAUDE\_CODE\_EXPERIMENTAL\_AGENT\_TEAMS} 环境变量。Team lead 分配任务给 teammates，每个 teammate 有独立上下文窗口和共享任务列表，可以直接互相通信。

建议从 3--5 个 teammates、每人 5--6 个任务开始。避免让两个 teammate 编辑同一文件（会导致覆写）。先从研究和审查任务开始练手。
\end{knowledgebox}

\subsection{一个 Claude 写代码，另一个 Claude 审查（Tip \#45）}

第一个 Claude 实现功能，第二个 Claude 从全新上下文出发像 Staff Engineer 一样审查。审查者不知道实现过程中的偷懒之处，会挑战每一个取巧。

同样的思路适用于 TDD：Session A 写测试，Session B 写代码通过测试。

\subsection{用 claude -p 做批量操作（Tip \#49）}

\begin{lstlisting}[caption={批量操作示例}, language=bash]
for file in src/*.ts; do
  claude -p "Update imports in $file" \
    --allowedTools Write,Read &
done
wait
\end{lstlisting}

\texttt{-{}-allowedTools} 限定每个文件允许的操作，\texttt{\&} 并行执行。适合格式转换、批量更新 import、独立文件的重复迁移。

\subsection{远程控制（Tip \#9）}

运行 \texttt{claude remote-control} 启动会话，然后从 \texttt{claude.ai/code} 或 iOS/Android 上的 Claude App 连接。会话在本地运行，手机/浏览器只是窗口。可以发消息、审批工具调用、随时监控进度。

\subsection{添加 Status Line（Tip \#18）}

Status line 是每个 Claude 回合后运行的 shell 脚本，在终端底部显示实时信息：当前目录、Git 分支、上下文使用量（按窗口占用率着色）。用 \texttt{/statusline} 快速设置。

\subsection{本章小结}

并行工作的层次：单任务用 Subagent 隔离调查成本 $\to$ 多分支用 Worktree 并行开发 $\to$ 跨会话用 Agent Teams 协作 $\to$ 批量文件用 \texttt{claude -p} 扇出处理。远程控制让你可以移动端监控长任务。

%% ============================================================
